% !PRINTED 246
\weberpage{246}
%<<CONT
\weq{5}{\Phi_{a,c,\partial} = \frac{\Psi\left[j\left(\frac{c+\partial\omega}{a}\right), j(\omega)\right]}{F'\left[j\left(\frac{c+\partial\omega}{a}\right)\right]}.}

Thus $\Phi_{a,c,\partial}$ belongs to the algebraic field which consists
of the rational functions of $j\left(\frac{c+\partial\omega}{a}\right)$
and $j(\omega)$, and we can apply the theorems of Bd. I, § 151. From
these it follows that the $\nu$ quantities $\Phi_{a,c,\partial}$ are
the roots of an equation of degree $\nu$ whose coefficients depend
rationally on $j(\omega)$. If the $\nu$ quantities
$\Phi_{a,c,\partial}$ are different from one another, then this
equation is \gesperrt{irreducible}. Such an equation we call a \gesperrt{transformation
equation of the first level belonging to the transformation degree
$n$}\footnote{In the first edition I called these equations
``\gesperrt{invariant transformation equations}''. This expression has
been objected to by \gesperrt{Klein} (Vorlesungen über ausgewählte
Kapitel der Zahlentheorie, autographiertes Heft, Göttingen 1897). I
adopt \gesperrt{Klein's} proposal, in that I now call these equations
``\gesperrt{transformation equations of the first level}'', without
entering here more closely into the justification of this
expression.}. Every other quantity of the field can be rationally
expressed by such a $\Phi_{a,c,\partial}$ and by $j(\omega)$.

If the functions $\Phi_{a,c,\partial}$ have the property of remaining
finite for every finite value of $\omega$ with positive imaginary
part, and hence for every finite $j(\omega)$, then the function
$\Psi[x,j(\omega)]$ in (4) is also a \gesperrt{polynomial}
in $j(\omega)$. Formula (5) could therefore fail only for
those \gesperrt{particular values} of $\omega$ for which two roots of the
invariant equation become equal to one another.

If the $\nu$ quantities $\Phi_{a,c,\partial}$ are not all different
from one another, then they split into groups of equal size, each group
consisting of equal quantities, and the equation of degree $\nu$ to be derived from (5) is
a power of an irreducible equation, which we shall occasionally also
designate as a transformation equation (Bd. I, § 151, 2).

If the quantities $\Phi_{a,c,\partial}$ are so chosen that they become
infinite for no finite $\omega$ with positive imaginary part, then
they remain finite for every finite $j(\omega)$, from which it follows
that the coefficients in the function of degree $\nu$
\weq{6}{\Pi(x-\Phi_{a,c,\partial})}
are \gesperrt{polynomials in $j(\omega)$}, i.e. the
$\Phi_{a,c,\partial}$ are \gesperrt{integral algebraic functions of
$j(\omega)$}.
